\section{The likelihood}\label{sec:likelihood}

The likelihood of a cohort is computed from the exact distribution of the first-passage time of the process in
eq.~\eqref{eq:sr}, not from a simulation or a fitted curve. This section gives the equation that is solved, its
discretisation, the observation model that turns survival and absorbed-mass curves into a log-likelihood, and the grid
rules that set the numerical resolution.

\paragraph{The Fokker--Planck problem.} Let $p(x,t)$ be the density of damage among individuals that are still alive
at age $t$ (with $A=1$, so $Y(t)=t$). It obeys the Fokker--Planck equation of eq.~\eqref{eq:sr} \citep{risken1996fokker},
\begin{equation}\label{eq:fp}
  \partial_t p = -\partial_x\!\Bigl[\Bigl(\eta\,t-\beta\frac{x}{x+\kappa}\Bigr)p\Bigr]+\eps\,\partial_x^2 p,
  \qquad p(x,0)=\delta_0(x),
\end{equation}
on $0<x<\xc=1$. The boundary at $0$ is reflecting (zero flux) and the boundary at $\xc$ is absorbing, $p(\xc,t)=0$. The
survival function is $S(t)=\int_0^{\xc}p\,\dd x$. Individuals die by crossing $\xc$, so the death density is the flux
absorbed at the boundary,
\begin{equation}\label{eq:flux}
  f(t)=-\dd S/\dd t=-\eps\,\partial_x p(\xc,t),\qquad F(t)=\int_0^t f(s)\,\dd s=1-S(t).
\end{equation}
Time enters through $\tau=t/L$, where $L$ is the Kaplan--Meier mean lifetime of the cohort
\citep{kaplan1958nonparametric}. In $\tau$ the problem has the rates $\rho_\eta=L^2\eta$, $\rho_\beta=L\beta$,
$\rho_\eps=L\eps$ and the unchanged $\kappa$, so that all quantities are of order one. An external hazard $h$ (Section~\ref{sec:model})
acts outside the equation: the observed survival is $S_h(t)=S(t)\,e^{-ht}$.

\paragraph{Discretisation.} The interval $(0,\xc)$ is divided into $N$ cells of width $\Delta y=1/N$ with cell values
$p_j$. The flux between cells $j$ and $j+1$ uses the Scharfetter--Gummel scheme \citep{scharfetter1969large}, which
stays positive and accurate when drift dominates diffusion:
\begin{equation}\label{eq:sg}
  J_{j+\frac12}=\frac{\rho_\eps}{\Delta y}\bigl[B(-P_j)\,p_j-B(P_j)\,p_{j+1}\bigr],\qquad
  B(z)=\frac{z}{e^{z}-1},\qquad P_j=\frac{(\rho_\eta\tau-r_j)\,\Delta y}{\rho_\eps},
\end{equation}
where $r_j=\rho_\beta y_j/(y_j+\kappa)$ is the removal rate at the cell centre. The flux is zero at the wall at $0$, and
the flux out of the last cell, $J_{N-\frac12}$, is the absorbed flux of eq.~\eqref{eq:flux}. Writing the resulting
tridiagonal operator as $A(\tau)$, time is advanced by backward Euler for the first steps, which damps the
$\delta$-function start, and then by the second-order backward differentiation formula BDF2 \citep{hairer1996solving},
\begin{equation}\label{eq:bdf2}
  \bigl(\tfrac32 I-\Delta\tau\,A(\tau_{n+1})\bigr)\,p^{n+1}=2p^{n}-\tfrac12 p^{n-1},\qquad p^{n+1}\leftarrow\max(p^{n+1},0),
\end{equation}
with the corresponding variable-step coefficients on a non-uniform mesh. The first \C{m-nstart} steps are backward Euler.
Clipping negative values to zero keeps the density a density; the mass it adds back is accounted for below. The march
stops when $S$ falls below \C{m-gtol} and $\tau>$\,\C{m-tau-guard}, since later deaths carry no probability, and $S$ is made
monotone afterwards. The absorbed mass $F$ is accumulated step by step from the flux $J_{N-\frac12}$ with the same
time-stepping weights, plus the mass restored by the clip. The reason is round-off. An early death has mass
$F(a+\delta)-F(a)$ that can be many orders of magnitude below one; computed as $S(a)-S(a+\delta)$ it is a difference of two
numbers near one and is lost in round-off, whereas every term accumulated into $F$ is small and positive and keeps full
relative precision. The interval mass is therefore taken from $F$ wherever $F$ is the smaller operand, and from $S$
otherwise; the two are equal in exact arithmetic. Gradients of the log-likelihood with respect to $\log\theta$ are
exact, obtained by a hand-written reverse sweep of the march.

\paragraph{Observation model.} A cohort has, for each individual $i$, a recorded age $a_i$, an event flag $d_i$ (1
death, 0 right-censored) and a recording resolution $\delta$. A death recorded at $a_i$ is scored on $[a_i,a_i+\delta)$.
Under census-based recording, an animal found dead at census $c$ was alive at the previous census, so it is entered
with $a_i=c-\delta$. The contribution of individual $i$ is
\begin{equation}\label{eq:term}
  \ell_i=d_i\log m_i+(1-d_i)\log S_h(a_i)-\log S_h(q_i),
\end{equation}
with the interval mass of a death and its external-hazard correction
\begin{equation}\label{eq:mass}
  m_i=e^{-h a_i}\Bigl[\bigl(F(a_i+\delta)-F(a_i)\bigr)+S(a_i+\delta)\bigl(1-e^{-h\delta}\bigr)\Bigr],
\end{equation}
which reduces to $F(a_i+\delta)-F(a_i)$ when $h=0$. For exact (continuous) death times $m_i$ is replaced by the density
$\bigl(f(a_i)+h\,S(a_i)\bigr)e^{-h a_i}$. The total is $\ell(\theta)=\sum_i\ell_i$, each term is floored at a log-probability of
\C{m-log-floor} so that one impossible death does not make the likelihood undefined, and $\ell=-\infty$ if the sum is
not finite. The quantity $q_i$ in eq.~\eqref{eq:term} is the entry age. It is zero
for individuals observed from birth; individuals first observed at age $q_i>0$ (left truncation, as in the group
comparisons of Section~\ref{sec:effects}) are conditioned on survival to $q_i$, which is the last term.

\emph{Right censoring.} A censored individual contributes $\log S_h(a_i)$, the probability of surviving to its last
observation.

\emph{Fit window $[\mathrm{lo},\mathrm{hi})$.} A window restricts the likelihood to the part of life in which the model is
trusted. Individuals recorded before $\mathrm{lo}$ are dropped, every remaining individual enters with $q_i=\mathrm{lo}$
(condition on $S_h(\mathrm{lo})$), and individuals alive or recorded as dying at or after $\mathrm{hi}$, including a death
whose recording interval crosses $\mathrm{hi}$, are censored at $\mathrm{hi}$. The window edges are snapped to the
recording grid.

\emph{Ties.} Individuals with the same recorded age have the same term; with interval scoring no tie-breaking rule is
needed. One tie is not harmless. If all censored individuals share one age $t_c$ equal to the last census, a death
recorded in the recording interval that contains $t_c$ may have occurred after $t_c$, while the censored are credited
with $S(t_c)$; scoring that death over its whole interval counts the sliver after $t_c$ twice and biases $\eta$ upward.
These deaths are therefore relabelled as censored at their own recorded age (Section~\ref{sec:estimation}).

\paragraph{Parameters that change with age.} The engine accepts parameter changes on age segments
$[t_s,t_{s+1})$, each with its own multipliers on $\eta$, $\beta$, $\eps$, $\kappa$, $\xc$ and $A$; this is how
Section~\ref{sec:effects} represents a treatment started at age $t_0$. Parameters are constant within a segment and the
density is carried across the edge unchanged. The production term uses the accumulated damage-production units $Y$, which grow
at the rate $A$ of the current segment (so $Y$ is continuous and only its slope changes at an edge), whereas removal, noise
and the time step are not rescaled. Segment edges are snapped to the nearest grid node; the step that starts at the edge is a
backward Euler restart. A change of the threshold moves the absorbing cell: raising it exposes empty cells, and lowering it
absorbs at once the mass above the new threshold, which is added to $F$ as a pulse of deaths at the edge. This is exact
on the grid; a threshold change is not equivalent to rescaling the other parameters except from birth
(Section~\ref{sec:model}).

\paragraph{Grid rules.} The numerical resolution is set from the data by rules that were calibrated on the benchmark
(Section~\ref{sec:val-numerics}); all constants are listed in Table~\ref{tab:switches} of the SI.
\begin{itemize}[leftmargin=1.5em,itemsep=1pt]
  \item \emph{Scale and horizon.} $L$ is the Kaplan--Meier mean. The horizon is a multiple of the last recorded age and the nominal
    step divides it into a fixed number of steps (SI Table~\ref{tab:switches}).
  \item \emph{Recording ceiling.} The step is at most $\delta/(3L)$, so that a recording interval spans at least three steps. Where this
    would fall below \C{m-dt-floor} the ceiling is held at that value and $N$ is doubled to \C{m-N-floor}.
  \item \emph{Graded early mesh.} The march starts with steps of \C{m-early-dt0} that grow by a factor \C{m-early-ratio} per step
    up to the nominal step. Deaths in the first recording intervals then keep a small grid error; with a uniform start they
    can be several nats off.
  \item \emph{The $N$ rule.} The number of cells is $N=\C{m-N0}\cdot 2^{k}$ for the smallest $k\ge0$ for which the predicted error of the
    interval-relevant deviances, $E_{900}\,q^{k}$, is at most \C{m-nr-target} nats, with a cap of \C{m-Nmax} cells. The error at $N=\C{m-N0}$ is
    predicted from the cohort size $n$ and the Kaplan--Meier coefficient of variation (CV) of the lifetime,
    $\log_{10}E_{900}=\C{nrule-A}+\C{nrule-B}\log_{10}n+(\C{nrule-C})\log_{10}\mathrm{CV}$, fitted on \C{nrule-n} cohorts, with the error
    contracting by $q=\C{m-nr-q}$ per doubling of $N$. The prediction comes from a grid study of \C{conv-n-measured} stored
    cohorts, chosen because they were numerically hard (Section~\ref{sec:val-numerics}, Figure~\ref{fig:convergence}). For
    each cohort, $E_{900}$ is the largest error of the deviance on $N=\C{nrule-N0}$ cells at the points that set the interval
    edges, measured against a Richardson limit from grids of up to $14400$ cells; the regression uses the \C{nrule-n} cohorts
    whose values are reproducible, and its residual standard deviation is \C{nrule-resid-sd} decades. The rule is a median
    model with no safety margin: the measured error exceeds the prediction for \C{conv-resid-pos}\,\% of these cohorts. The
    contraction $q$ is a chosen value. It is slower than the measured median on non-human cohorts
    (\C{conv-r-med-nonhuman}) but faster than on human cohorts (\C{conv-r-med-human}), so the rule under-refines large human
    life tables.
  \item \emph{The step rule for sharp cohorts (\code{DT_RULE}).} When lifetimes are tightly clustered, the time error dominates
    and the step is capped at $\mathrm{CV}/\C{m-dtr-k}$. The rule applies only where the CV, corrected for the recording
    resolution, $\mathrm{CV}_c=\sqrt{\max(0,\mathrm{CV}^2-(\delta/L)^2/12)}$, is below \C{m-dtr-cvmax}. The correction removes the
    spread that recording adds: a death known only to lie in an interval of length $\delta$ is, within it, uniform, and a
    uniform distribution over an interval of length $\delta$ has variance $\delta^2/12$ (in units of $L$, $(\delta/L)^2/12$).
    Without the rule, BDF2 overshoot and the
    positivity clip set $S$ to zero past the mean and every later death falls on the numerical floor.
\end{itemize}
Cost is proportional to $N$ times the number of steps, so the rules for large cohorts and sharp cohorts increase run time
(Section~\ref{sec:estimation}, Section~\ref{sec:limitations}).
\paragraph{Heterogeneity.} A population with a spread of one parameter (for example the threshold) is a mixture over
Gauss--Hermite nodes of that parameter: the survival, absorbed-mass and density curves are averaged over nodes, and the
observation terms of eq.~\eqref{eq:term} are applied once to the mixture, so that conditioning on entry divides the
mixture rather than mixing conditional terms. This is available in the engine but is not part of the released calibrators
and is not validated for calibration (Section~\ref{sec:limitations}).
